\section{二元关系的逆}

\begin{lemma}{二元关系的逆的存在性}{}
	设$G$是二元关系,则公式$\left(y,x\right)\in G$有图像.
\end{lemma}

\begin{definition}{二元关系的逆}{}
	设$G$是二元关系.所有满足$\left(y,x\right)\in G$的有序对$\left(x,y\right)$组成的二元关系称为$G$的逆关系,记作$G^{-1}$(或$G^{\op}$).
\end{definition}

\begin{definition}{集合在二元关系下的原像}{}
	设$G$是二元关系.对任一集合$X$,集合$G^{-1}\left(X\right)$称为$X$在$G$下的原像.
\end{definition}

\begin{proposition}{逆关系的基本性质}{}
	设$G$是二元关系.
	\begin{enumerate}
		\item $\left(G^{-1}\right)^{-1}=G$.
		\item $\pr_{1}\left(G^{-1}\right)=\pr_{2}\left(G\right),\pr_{2}\left(G^{-1}\right)=\pr_{1}\left(G\right)$.
	\end{enumerate}
\end{proposition}

\begin{proposition}{笛卡尔积的逆}{}
	设$X,Y$是集合,则$\left(X\times Y\right)^{-1}=Y\times X$.
\end{proposition}

\begin{definition}{对称二元关系}{}
	设$G$是二元关系.若$G^{-1}=G$,则称$G$是对称的.
\end{definition}

\begin{definition}{逆对应}{}
	设$\Gamma\coloneq\left(G,A,B\right)$是对应.我们称$\left(G^{-1},B,A\right)$是$\Gamma$的逆对应,记作$\Gamma^{-1}$.
\end{definition}

\begin{definition}{集合在对应下的原像}{}
	设$\Gamma\coloneq\left(G,A,B\right)$是对应.对$B$的任意子集$Y$,$\Gamma^{-1}\left(Y\right)$称为$Y$在$\Gamma$下的原像.
\end{definition}

\begin{proposition}{逆对应的对合性}{}
	设$\Gamma\coloneq\left(G,A,B\right)$是对应,则$\left(\Gamma^{-1}\right)^{-1}=\Gamma$.
\end{proposition}















